% Ensure included PDFs (banners) with newer versions embed without warnings
\pdfminorversion=7
\documentclass[aspectratio=169]{beamer}

\makeatletter
% Allow compiling from either Slides/ or repository root.
\def\input@path{{Template/}{Slides/Template/}}
\makeatother
\usepackage{graphicx}
% Make graphics (including banner PDFs) resolvable without TEXINPUTS
\graphicspath{{Template/}{Template/banners/}{Slides/Template/}{Slides/Template/banners/}{../Common_Images/}{Common_Images/}}

%================================================================%
% Theme and Package Setup
%================================================================%
\usetheme{ucl}
\setbeamercolor{banner}{bg=darkpurple}

% Navigation and Footer
\setbeamertemplate{navigation symbols}{\vspace{-2ex}}
\setbeamertemplate{footline}[author title date]
\setbeamertemplate{slide counter}[framenumber/totalframenumber]

\usepackage[utf8]{inputenc}
\usepackage[british]{babel}
\usepackage{amsmath, amssymb, amsthm}
\usepackage{tikz}
\usepackage{pgfplots}
\pgfplotsset{compat=1.17}
\usetikzlibrary{calc, positioning, arrows.meta, shapes.geometric, trees, backgrounds, shapes.misc, graphs, quotes, shadows, fit, decorations.pathreplacing} 
\usefonttheme{professionalfonts}
\usepackage{eulervm}
\usepackage{listings}
\usepackage{xcolor}
\usepackage{algorithm}
\usepackage{algpseudocode}
\usepackage[square,authoryear]{natbib}

% Define custom colours
\makeatletter
\@ifundefined{color@stone}{%
    \definecolor{stone}{gray}{0.95}%
}{}
\makeatother
\definecolor{darkgreen}{rgb}{0.0, 0.5, 0.0}
\definecolor{uclblue}{cmyk}{1,0.24,0,0.64}

\setbeamercovered{transparent}

% Define theoremblock
\makeatletter
\newenvironment<>{theoremblock}[1]{%
    \begin{block}#2{#1}%
}{\end{block}}
\makeatother

% Section divider slides
\AtBeginSection[]{
  \begin{frame}
    \frametitle{\textbf{\Large\insertsectionhead}}
    \begin{center}
        \huge\textbf{\insertsectionhead}
    \end{center}
  \end{frame}
}

% Operator definitions
\DeclareMathOperator*{\argmin}{arg\,min}
\DeclareMathOperator{\prox}{prox}
\DeclareMathOperator{\proj}{proj}
\DeclareMathOperator{\divg}{div}

\title{Modular Deep Learning and Generative AI}
\subtitle{From Plug-and-Play to Diffusion Models}
\author[Martin Benning (University College London)]{Martin Benning}
\date[CCMI]{CCMI -- Deep Learning in Theory and Practice \\[1cm] Wednesday, 20 May 2026}
\institute[]{University College London}

\begin{document}

% --- Title Frame ---
\begin{frame}
  \titlepage
\end{frame}

% --- Overview Frame ---
\begin{frame}{Lecture Overview}
  \begin{block}{Objectives for Today}
    \begin{itemize}
        \item Overcome the rigidity of end-to-end and unrolled networks via \textbf{Plug-and-Play (PnP)}.
        \item Critically evaluate explicitly defined regularisation losses using \textbf{Regularisation by Denoising (RED)} and the \textbf{Jacobian Critique}.
        \item Shift from deterministic optimisation (MAP) to \textbf{Probabilistic Generative Models} (VAEs, GANs, Normalising Flows).
        \item Connect MSE-trained denoisers to probability distributions using \textbf{Tweedie's Formula}.
        \item Utilise \textbf{Langevin Dynamics} and \textbf{Diffusion Posterior Sampling (DPS)} to solve inverse problems by sampling from the true posterior.
    \end{itemize}
  \end{block}
\end{frame}

\begin{frame}{Lecture Overview}
  \tableofcontents
\end{frame}

%================================================================%
\section{Modular Deep Learning via Plug-and-Play}
%================================================================%

\begin{frame}{The End-to-End Problem (Recap)}
    On Tuesday, we derived highly powerful unrolled architectures like the Learned Primal-Dual (LPD) network. While these models represent the state-of-the-art for specific tasks, they suffer from a fundamental flaw: \textbf{Rigidity}.\pause
    
    \vspace{0.5em}
    By unrolling the iterations, the network incorporates the forward operator $A$ directly into its computational graph during training.\pause
    
    \begin{alertblock}{The Retraining Bottleneck}
        If the physics change even slightly -- for instance, altering the resolution, varying the sampling mask in MRI, or switching to a different sensor geometry -- the unrolled network must be entirely retrained.
    \end{alertblock}
    
    \vspace{0.5em}
    Can we decouple the neural network from the physical operator $A$, creating a \emph{modular} system?
\end{frame}

\begin{frame}{Variable Splitting as a DL Framework}
    To achieve modularity, we revisit classical variable splitting algorithms. We previously discussed PDHG, but today we introduce a simpler method known as \textbf{Half-Quadratic Splitting (HQS)}.
    
    Consider solving our standard variational problem via HQS, i.e.
    \begin{align*}
        \min_{x, v} \left\{ \frac{1}{2}\|Ax - y^\delta\|_2^2 + \alpha \mathcal{J}(v) + \frac{\mu}{2} \|x - v\|_2^2 \right\} \, .
    \end{align*}\pause
    Here, we introduce an auxiliary variable $v$ and penalise its distance to the image $x$ with a parameter $\mu > 0$. As $\mu \to \infty$, $x$ converges to $v$, we solve the original problem.
\end{frame}

\begin{frame}{Alternating Minimisation in HQS}
    Similar to ADMM, we solve the HQS objective via alternating minimisation yielding
    \begin{align*}
        x^{k+1} &= \argmin_{x} \left\{ \frac{1}{2}\|Ax - y^\delta\|_2^2 + \frac{\mu}{2} \|x - v^k\|_2^2 \right\} \, , \\
        v^{k+1} &= \argmin_{v} \left\{ \frac{\mu}{2} \|x^{k+1} - v\|_2^2 + \alpha \mathcal{J}(v) \right\} \, .
    \end{align*}
    
    \pause
    \vspace{0.5em}
    We have successfully decoupled the forward operator $A$ from the potentially non-linear, non-smooth regulariser $\mathcal{J}$.
\end{frame}

\begin{frame}{Analysing the Sub-Problems}
    Notice how HQS elegantly decouples the problem:
    \begin{enumerate}
        \item \textbf{Data Module ($x$-update):} For Gau\ss ian noise, this requires solving a simple linear system involving only the forward operator $A$ and the measurements $y^\delta$, which reads
        \begin{align*}
            x^{k+1} = (A^\top A + \mu I)^{-1} (A^\top y^\delta + \mu v^k) \, .
        \end{align*}
        \item \textbf{Prior Module ($v$-update):} By dividing by $\mu$, we recognise the exact definition of the proximal operator for the function $\frac{\alpha}{\mu} \mathcal{J}$ given by
        \begin{align*}
            v^{k+1} = \operatorname{prox}_{\frac{\alpha}{\mu} \mathcal{J}} (x^{k+1}) = \argmin_{v} \left\{ \frac{1}{2} \|v - x^{k+1}\|_2^2 + \frac{\alpha}{\mu} \mathcal{J}(v) \right\} \, .
        \end{align*}
    \end{enumerate}
\end{frame}

\begin{frame}{The Denoising Insight}
    What exactly is the prior module doing?
    
    \vspace{0.5em}
    The proximal step asks for a vector $v$ that is close to the noisy intermediate input $x^{k+1}$, but strongly regularised by $\mathcal{J}(v)$. 
    
    \vspace{0.5em}
    If we interpret $\mathcal{J}(v)$ as a negative log-prior (i.e., $\mathcal{J}(v) = -\log p(v)$), then this step is mathematically strictly equivalent to a \textbf{Maximum A-Posteriori (MAP) Gau\ss ian Denoising} problem!
\end{frame}

\begin{frame}{The Plug-and-Play (PnP) Heuristic}
    The core insight of the \textbf{Plug-and-Play} heuristic is that we do not need to manually define a mathematical prior $\mathcal{J}$ and derive its proximal operator \citep{venkatakrishnan2013plug,zhang2021plug}.
    
    \begin{theoremblock}{The PnP Strategy}
        We simply replace the proximal operator with a general, state-of-the-art neural network denoiser $D_\sigma(\cdot)$ that has been pre-trained on vast datasets of natural images:\vspace{-0.15cm}
        \begin{align*}
            v^{k+1} = D_{\sigma_k}(x^{k+1}) \, .
        \end{align*}
        
        \vspace{-0.15cm}
    \end{theoremblock}\pause
    
    By ``plugging in'' highly advanced models like BM3D \citep{dabov2007image}, DnCNN \citep{zhang2017beyond}, UNet \citep{ronneberger2015u}, or DRUNet, we leverage deep learning purely as a universal spatial prior, maintaining total independence from the specific measurement physics $A$ \citep{milanfar2025denoising,tan2025solving}.
\end{frame}

\begin{frame}{PnP-HQS Algorithm}
    Combining HQS with the PnP heuristic gives us our first practical algorithm.
    
    \begin{algorithm}[H]
    \caption{Plug-and-Play HQS}\label{alg:pnphqs}
    \begin{algorithmic}[1]
    \State \textbf{Input:} Data $y^\delta$, Operator $A$, Denoiser $D_\sigma$, Params $\mu, \alpha$
    \State \textbf{Initialise:} $v^0 = y^\delta$, $x^0 = A^\top y^\delta$
    \For{$k = 0, \dots, N-1$}
        \State \textbf{1. Inversion:} \ \ \ $x^{k+1} = (A^\top A + \mu I)^{-1} (A^\top y^\delta + \mu v^k)$
        \State \textbf{2. Denoising:} \ \ $\sigma_k = \sqrt{\alpha / \mu}$
        \State \ \ \ \ \ \ \ \ \ \ \ \ \ \ \ \ \ \ \ \ \ \ $v^{k+1} = D_{\sigma_k}(x^{k+1})$
    \EndFor
    \State \textbf{Return:} $x^{N}$
    \end{algorithmic}
    \end{algorithm}
\end{frame}

\begin{frame}{From HQS to ADMM}
    While HQS is incredibly simple, the Alternating Direction Method of Multipliers (ADMM) handles constraints properly and avoids the pitfalls of penalty methods (where $\mu \to \infty$ is needed to strictly enforce $x=v$).\pause
    
    \vspace{0.5em}
    The Augmented Lagrangian for ADMM, where we substitute the original Lagrange multiplier directly with a scaled dual variable $w$, reads
    \begin{align*}
        \mathcal{L}_\mu(x, v; w) = \frac{1}{2} \| Ax - y^\delta \|_2^2 + \alpha \mathcal{J}(v) + \mu \langle w, x - v \rangle + \frac{\mu}{2} \| x - v \|_2^2 \, .
    \end{align*}
\end{frame}

\begin{frame}{Conventional ADMM Algorithm}
    The classical ADMM alternatingly updates the variables of this augmented Lagrangian:
    \vspace{-0.15cm}
    
    \begin{algorithm}[H]
    \caption{Conventional ADMM}\label{alg:admm}
    \begin{algorithmic}[1]
    \State \textbf{Input:} Data $y^\delta$, Operator $A$, Regulariser $\mathcal{J}$, Parameters $\mu, \alpha$
    \State \textbf{Initialise:} $x^0 = A^\top y^\delta$, $v^0 = y^\delta$, $w^0 = 0$
    \For{$k = 0, 1, \dots$}
        \State \textbf{1. Data Consistency:} $x^{k+1} = (A^\top A + \mu I)^{-1} \left( A^\top y^\delta + \mu (v^k - w^k) \right)$
        \State \textbf{2. Proximal Prior:} $v^{k+1} = \operatorname{prox}_{\frac{\alpha}{\mu}\mathcal{J}}(x^{k+1} + w^k)$
        \State \textbf{3. Dual Update:} $w^{k+1} = w^k + (x^{k+1} - v^{k+1})$
    \EndFor
    \end{algorithmic}
    \end{algorithm}
\end{frame}

\begin{frame}{PnP-ADMM Algorithm}
    Replacing the proximal map in the classical ADMM updates with $D_\sigma$ provides the highly robust \textbf{PnP-ADMM} algorithm:
    \vspace{-0.15cm}
    
    \begin{algorithm}[H]
    \caption{Plug-and-Play ADMM}\label{alg:pnp_admm}
    \begin{algorithmic}[1]
    \State \textbf{Input:} Data $y^\delta$, Operator $A$, Pre-trained Denoiser $D_\sigma$, Penalty $\mu$
    \State \textbf{Initialise:} $x^0 = A^\top y^\delta$, $v^0 = y^\delta$, $w^0 = 0$
    \For{$k = 0, 1, \dots$}
        \State \textbf{1. Data Consistency:} $x^{k+1} = (A^\top A + \mu I)^{-1} \left( A^\top y^\delta + \mu (v^k - w^k) \right)$
        \State \textcolor{red}{\textbf{2. Neural Denoising (PnP change):} $v^{k+1} = D_{\sigma}(x^{k+1} + w^k)$}
        \State \textbf{3. Dual Update:} $w^{k+1} = w^k + (x^{k+1} - v^{k+1})$
    \EndFor
    \end{algorithmic}
    \end{algorithm}
    
    \vspace{-0.15cm}
    \textbf{Note:} The CNN denoiser cleans the sum of the primal and dual variables, effectively projecting the output back onto the manifold of natural images.
\end{frame}

\begin{frame}{Practical Parameter Tuning}
    A crucial aspect of PnP is the relationship between the regularisation parameter $\alpha$ and the internal noise level of the denoiser $\sigma$.
    
    \vspace{0.5em}
    In theory, for the MAP interpretation to hold, we require $\sigma_k = \sqrt{\alpha/\mu}$.\pause
    
    \begin{block}{The Continuation Method}
        In practice, when using deep denoisers (like DRUNet), it is highly beneficial to \textbf{decrease $\sigma_k$} dynamically over iterations. Starting with a large $\sigma$ removes heavy artifacts and poor initialisations, while a small $\sigma$ in later iterations refines fine details without washing away contrast.
    \end{block}
\end{frame}

\begin{frame}{Stability: The Convergence Problem}
    While the empirical success of PnP is undeniable, establishing rigorous mathematical guarantees for convergence remains challenging.
    
    \vspace{0.5em}
    Classical splitting proofs rely heavily on the assumption that the regularisation term $\mathcal{J}$ is a proper, closed, and convex function, allowing us to form a non-expansive sequence that monotonously decreases an Augmented Lagrangian.
    
    \begin{alertblock}{Breaking the Assumption}
        When we replace the proximal step with an off-the-shelf CNN denoiser $D_\sigma(\cdot)$, we fundamentally break this assumption. Most state-of-the-art denoisers do not correspond to the proximal operator of \emph{any} underlying regularisation functional.
    \end{alertblock}
\end{frame}

\begin{frame}{1. The Optimisation Perspective}
    To preserve classical convergence proofs, one must ensure that the chosen denoiser $D_\sigma$ mathematically behaves like a proximal operator. 
    
    \vspace{0.5em}
    By Moreau's Theorem, a function $D_\sigma$ is the proximal operator of some closed, proper, convex function if and only if \citep{moreau1965proximite,combettes2011proximal,Bauschke2011}:
    \begin{enumerate}
        \item \textbf{Non-expansiveness:} $D_\sigma$ must be Lipschitz continuous with a Lipschitz constant $L \le 1$.
        \item \textbf{Symmetric Jacobian:} The Jacobian matrix $\nabla D_\sigma(x)$ must be perfectly symmetric for all $x$, meaning $D_\sigma$ forms a conservative vector field.
    \end{enumerate}
    
    \vspace{0.5em}
    Because standard denoisers violate the symmetric Jacobian property, researchers must design tailored neural network architectures to strictly enforce this.
\end{frame}

\begin{frame}{Stability in Deep Learning Feedback Loops}
    \textbf{2. The Fixed-Point Perspective:} If we abandon the requirement that an underlying objective function exists, we must analyse PnP as a sequence of operator compositions seeking a state of "consensus equilibrium".
    \vspace{0.5em}
    
    The entire PnP-HQS algorithm can be written as a single operator iteration $v^{k+1} = T(v^k)$, where
    \begin{align*}
        T = D_\sigma \circ \operatorname{prox}_{G/\mu} \, .
    \end{align*}
    
    \begin{itemize}
        \item \textbf{Strict Contraction ($L_D < 1$):} By the Banach Fixed-Point Theorem, PnP-HQS converges uniquely regardless of initialisation.
        \item \textbf{Non-Expansive ($L_D \le 1$):} Convergence can be guaranteed by employing the Krasnosel'skii-Mann (KM) theorem, applying an averaging/relaxation step.
    \end{itemize}
    We enforce $L_D \le 1$ during CNN training via \textbf{Spectral Normalisation}.
\end{frame}

%================================================================%
\section{Regularisation by Denoising (RED)}
%================================================================%

\begin{frame}{Regularisation by Denoising (RED)}
    PnP is an empirical success, but it breaks the link to variational optimisation. Because we replaced $\operatorname{prox}_{\mathcal{J}}$ arbitrarily, we no longer know what global energy landscape the algorithm is descending.
    
    \vspace{0.5em}
    \textbf{Regularisation by Denoising (RED)} attempts to construct an explicit regularisation functional directly from the chosen denoiser yielding
    \begin{align*}
        \rho(x) = \frac{1}{2} \langle x, x - D(x) \rangle \, .
    \end{align*}
    
    \textbf{Intuition:} The residual $x - D(x)$ isolates high-frequency noise. If $x$ is a clean image, $D(x) \approx x$, so the energy $\rho(x) \approx 0$. If $x$ is noisy, the penalty is large.
\end{frame}

\begin{frame}{Deriving the RED Gradient}
    To use $\rho(x)$ in a standard variational framework, we formulate the overall objective function yielding
    \begin{align*}
        E(x) = \mathcal{D}_{y^\delta}(Ax) + \alpha \rho(x) \, .
    \end{align*}
    To minimise this, we need the gradient $\nabla \rho(x)$. By the product rule, we obtain
    \begin{align*}
        \nabla \rho(x) &= \frac{1}{2} \nabla_x \left( \langle x, x \rangle - \langle x, D(x) \rangle \right) \, , \\
        &= \frac{1}{2} \left( 2x - D(x) - [\nabla D(x)]^\top x \right) \, , \\
        &= \frac{1}{2} \left( x - D(x) \right) + \frac{1}{2} \left( x - [\nabla D(x)]^\top x \right) \, ,
    \end{align*}
    where $\nabla D(x)$ is the Jacobian matrix of the denoiser. 
\end{frame}

\begin{frame}{The Gradient Simplification Claim}
    \citet{romano2017little} claimed that under specific conditions, the RED gradient simplifies directly to the denoising residual, yielding
    \begin{align*}
        \nabla \rho(x) = x - D(x) \, .
    \end{align*}
    For this to hold, we strictly require that $[\nabla D(x)]^\top x = D(x)$. The authors justify this via two strong mathematical assumptions about the denoiser:
    \begin{enumerate}
        \item \textbf{Jacobian Symmetry:} $[\nabla D(x)]^\top = \nabla D(x)$.
        \item \textbf{Local Homogeneity:} $D(cx) = c D(x)$ for scalar $c \approx 1$, implying the directional derivative condition $\nabla D(x) x = D(x)$.
    \end{enumerate}
    If true, one can simply deploy gradient descent: $x^{k+1} = x^k - \tau \left( A^\top(Ax^k - y^\delta) + \alpha (x^k - D(x^k)) \right)$.
\end{frame}

\begin{frame}{The Jacobian Critique}
    The RED framework offers an attractive narrative, but a detailed analysis by \citet{reehorst2018regularization} revealed a fundamental flaw.
    
    \vspace{0.5em}
    For an iterative algorithm updating via $x^{k+1} = x^k - \tau V(x^k)$ to minimise a scalar objective, the vector field $V(x)$ must be \textbf{conservative} (it must actually be the gradient of some scalar potential).
    
    \vspace{0.5em}
    A vector calculus theorem states that a continuously differentiable vector field $V$ is conservative if and only if its Jacobian matrix is perfectly symmetric. 
\end{frame}

\begin{frame}{Checking the Symmetry Condition}
    For the RED vector field $V(x) = x - D(x)$, the Jacobian is given by
    \begin{align*}
        \nabla V(x) = I - \nabla D(x) \, .
    \end{align*}
    Therefore, for the RED residual to be the gradient of \emph{any} scalar function, the denoiser's Jacobian $\nabla D(x)$ must be strictly symmetric, i.e.,
    \begin{align*}
        \frac{\partial [D(x)]_i}{\partial x_j} = \frac{\partial [D(x)]_j}{\partial x_i} \quad \text{for all } i, j \, .
    \end{align*}
\end{frame}

\begin{frame}{Reality Check: Do CNN Denoisers have Symmetric Jacobians?}
    \citet{reehorst2018regularization} conducted rigorous tests to determine if modern denoisers exhibit symmetric Jacobians.
    
    \begin{itemize}
        \item \textbf{Classical Denoisers:} Algorithms like Non-Local Means \citep{buades2005non} and BM3D \citep{dabov2007image} rely on complex, data-dependent block-matching mechanisms. Consequently, their Jacobians are highly asymmetric.
        \item \textbf{Deep Learning Denoisers (CNNs):} Standard feed-forward networks (e.g., DnCNN \citep{zhang2017beyond}, UNet \citep{ronneberger2015u}) involve interleaved convolutions and non-linear activations (like ReLU). Their Jacobians are entirely asymmetric.
    \end{itemize}
    
    \vspace{0.5em}
    \textbf{The Consequence:} The RED vector field has non-zero curl. It is not conservative, and therefore \textbf{minimises no objective function}. The elegant RED objective is a phantom function!
\end{frame}

\begin{frame}{Consensus Equilibrium}
    If standard RED algorithms (and PnP methods) are not minimising a cost function, how do we mathematically characterise their stationary points?
    
    \vspace{0.5em}
    Instead of viewing these methods through the lens of optimisation, we interpret them through the lens of \textbf{Consensus Equilibrium}. The algorithms are root-finding methods attempting to solve the system of non-linear equations yielding
    \begin{align*}
        A^\top(Ax^* - y^\delta) + \alpha (x^* - D(x^*)) = 0 \, .
    \end{align*}
    At a stationary point $x^*$, the algorithm achieves an equilibrium between two competing vector fields. The forward model gradient pulls towards data consistency, while the denoising residual pulls towards the manifold of natural images.
\end{frame}

\begin{frame}{Fixed-Point RED (FP-RED)}
    While a naive gradient descent approach can be used to seek the consensus equilibrium, \citet{romano2017little} proposed more efficient solvers.
    
    \vspace{0.5em}
    By rearranging the RED stationary condition to isolate one instance of $x^*$, we obtain a fixed-point equation yielding
    \begin{align*}
        (A^\top A + \alpha I) x^* = A^\top y^\delta + \alpha D(x^*) \, .
    \end{align*}
    This naturally suggests an iterative fixed-point scheme. Notice the striking structural similarity between this and the inversion step of PnP-HQS; both require solving a Tikhonov-regularised linear system at each step.
\end{frame}

\begin{frame}{The FP-RED Algorithm}
    We evaluate the denoiser at the current iterate $x^k$ and solve the resulting linear system for the next iterate $x^{k+1}$.
    
    \begin{algorithm}[H]
    \caption{Fixed-Point RED (FP-RED)}
    \begin{algorithmic}[1]
    \State \textbf{Input:} Data $y^\delta$, Operator $A$, Denoiser $D$, Parameter $\alpha$
    \State \textbf{Initialise:} $x^0 = A^\ast y^\delta$
    \For{$k = 0, \dots, N-1$}
        \State $x^{k+1} = (A^\top A + \alpha I)^{-1} \left( A^\top y^\delta + \alpha D(x^k) \right)$
    \EndFor
    \State \textbf{Return:} $x^{N}$
    \end{algorithmic}
    \end{algorithm}
\end{frame}

\begin{frame}{Convergence via Monotone Operator Theory}
    Because we lack a valid global objective function, we must evaluate FP-RED as a root-finding sequence solving a consensus equilibrium.
    
    \vspace{0.5em}
    The linear inversion operator $(A^\top A + \alpha I)^{-1}$ is firmly non-expansive, acting as a dampening factor. The convergence of the sequence $x^{k+1} = T_{FP}(x^k)$ relies entirely on the Lipschitz continuity of the denoiser $D(x)$:
    
    \begin{itemize}
        \item If the denoiser is a \textbf{strict contraction} ($L_D < 1$), FP-RED converges to a unique consensus equilibrium regardless of initialisation (Banach Fixed-Point Theorem).
        \item If the denoiser is \textbf{non-expansive} ($L_D \le 1$), applying a relaxation step guarantees convergence via the Krasnosel'skii-Mann (KM) theorem.
    \end{itemize}
\end{frame}

\begin{frame}{Gradient Step Denoisers (GSD)}
    To bridge this gap between rigorous theory and deep learning practice, we must structuralise the network to guarantee symmetry.
    
    \vspace{0.5em}
    Instead of defining $D(x)$ as a standard CNN, we define it explicitly as a gradient step on a learned convex potential $g_\theta$ yielding
    \begin{align*}
        D(x) = x - \nabla g_\theta(x) \, .
    \end{align*}
    We can parameterise $g_\theta$ using an \textbf{Input Convex Neural Network (ICNN)}.
\end{frame}

\begin{frame}{Properties of GSD}
    By structuring the denoiser as a gradient step, we fundamentally resolve the theoretical issues:
    
    \begin{itemize}
        \item \textbf{Symmetry Guaranteed:} The Jacobian $\nabla D = I - \nabla^2 g_\theta$ is defined by a Hessian matrix. Hessian matrices are inherently symmetric.
        \item \textbf{Conservative Field:} The denoiser actually represents the gradient of the potential $\phi(x) = \frac{1}{2}\|x\|^2 - g_\theta(x)$.
        \item \textbf{True Variational Optimisation:} Algorithms deploying a GSD are mathematically guaranteed to minimise a global energy landscape.
    \end{itemize}
\end{frame}

%================================================================%
\section{Deep Generative Modelling Frameworks}
%================================================================%

\begin{frame}{Shifting Paradigms: From Modes to Distributions}
    So far, our algorithms (unrolling, PnP, RED) have been completely deterministic. Given noisy data $y^\delta$, they climb a probability landscape to find a single, "best" estimate (the MAP estimate or consensus equilibrium).
    
    \vspace{0.5em}
    However, in highly ambiguous inverse problems, the posterior distribution $p(x | y^\delta)$ might possess multiple valid modes. Finding one mode ignores the vast uncertainty of the problem.
    
    \vspace{0.5em}
    \begin{theoremblock}{The Bayesian Shift}
        Instead of finding the single best output (MAP optimisation), we must model the entire data distribution $p_{\text{data}}(x)$ to evaluate expectations and sample diverse valid solutions from the posterior.
    \end{theoremblock}
\end{frame}

\begin{frame}{The Bayesian Framework \& KL Divergence}
    Our goal is to find a parameterised density function $p_\theta(x)$ that closely matches the true underlying data distribution $p_{\text{data}}(x)$.
    
    \pause
    \vspace{0.5em}
    We formulate this as an optimisation problem minimising the \textbf{Kullback-Leibler (KL) Divergence} between the two distributions yielding
    \begin{align*}
        D_{KL}(p_{\text{data}} \| p_\theta) = \mathbb{E}_{x \sim p_{\text{data}}}\left[ \log \frac{p_{\text{data}}(x)}{p_\theta(x)} \right] \, .
    \end{align*}
    
    \pause
    The KL-divergence is non-negative and strictly equals zero if and only if the distributions match perfectly almost everywhere.
\end{frame}

\begin{frame}{KL Divergence as a Bregman Distance (1/2)}
    On Monday and Tuesday, we introduced the \textbf{Bregman distance} for a convex function $J$ yielding
    \begin{align*}
        D_J(p, q) = J(p) - J(q) - \langle \nabla J(q), p - q \rangle \, .
    \end{align*}
    Can we frame the KL divergence $D_{KL}(p_{\text{data}} \| p_\theta)$ within this exact same mathematical framework?
    
    \pause
    \vspace{0.5em}
    If we choose $J(p) = \int p(x) \log p(x) dx$ (the negative Shannon entropy), the first variation (functional derivative) is $\nabla J(q) = \log q(x) + 1$. Substituting this into the Bregman distance definition provides
    \begin{align*}
        D_J(p, q) &= \int p(x) \log p(x) dx - \int q(x) \log q(x) dx \\
        &\quad - \int \left( \log q(x) + 1 \right) \left( p(x) - q(x) \right) dx \, .
    \end{align*}
\end{frame}

\begin{frame}{KL Divergence as a Bregman Distance (2/2)}
    We expand the inner product integral to observe
    \begin{align*}
        D_J(p, q) &= \int p(x) \log p(x) dx - \int q(x) \log q(x) dx \\
        &\quad - \int p(x) \log q(x) dx + \int q(x) \log q(x) dx \\
        &\quad - \int p(x) dx + \int q(x) dx \, .
    \end{align*}
    
    \pause
    Because $p$ and $q$ are valid probability density functions, we know $\int p(x) dx = 1$ and $\int q(x) dx = 1$. The constant terms and the $q \log q$ terms perfectly cancel yielding
    \begin{align*}
        D_J(p, q) = \int p(x) \log p(x) dx - \int p(x) \log q(x) dx = \mathbb{E}_{x \sim p}\left[ \log \frac{p(x)}{q(x)} \right] \, .
    \end{align*}
    
    \pause
    This strictly equals $D_{KL}(p \| q)$. %Thus, Generative AI training is intrinsically a Bregman minimisation problem.
\end{frame}

\begin{frame}{Maximum Likelihood Training}
    We train our generative model by minimising the KL-divergence. Expanding the expectation yields
    \begin{align*}
        \argmin_\theta D_{KL}(p_{\text{data}} \| p_\theta) &= \argmin_\theta \left( \mathbb{E}_{x \sim p_{\text{data}}}[\log p_{\text{data}}(x)] - \mathbb{E}_{x \sim p_{\text{data}}}[\log p_\theta(x)] \right) \, .
    \end{align*}
    
    \pause
    Since the first term is independent of the model parameters $\theta$, we can drop it. Approximating the expectation over our dataset $\mathcal{D} = \{x_i\}_{i=1}^N$ provides
    \begin{align*}
        \argmin_\theta D_{KL}(p_{\text{data}} \| p_\theta) \approx \argmin_\theta -\frac{1}{N} \sum_{i=1}^N \log p_\theta(x_i) \, .
    \end{align*}
    
    \pause
    This effectively results in \textbf{Maximum Likelihood Training}. We choose $\theta$ such that $p_\theta$ assigns the highest possible likelihood to the dataset.
\end{frame}

\begin{frame}{A Taxonomy of Generative AI}
    To implement this, we require neural network architectures that guarantee $p_\theta(x) \ge 0$ and $\int p_\theta(x) dx = 1$. We broadly divide these into Explicit and Implicit density models \citep{goodfellow2016}.
    
    \pause
    \vspace{0.5em}
    \begin{columns}
        \column{0.5\textwidth}
        \textbf{Explicit Density Models} \\
        Directly parameterise the probability density function $p_\theta(x)$.
        \begin{itemize}
            \item \textbf{Energy-Based Models (EBMs):} Unnormalised densities requiring a partition function $Z_\theta$.
            \item \textbf{Normalising Flows:} Exact likelihood via invertible diffeomorphisms and the change-of-variables formula.
        \end{itemize}
        
        \column{0.5\textwidth}
        \textbf{Implicit Density Models} \\
        Learn a generator $G_\theta(z)$ that maps a simple latent distribution ($z \sim \mathcal{N}(0, I)$) to complex data.
        \begin{itemize}
            \item \textbf{VAEs:} Maximise the Evidence Lower Bound (ELBO) using an encoder-decoder.
            \item \textbf{GANs:} Adversarial min-max game bypassing likelihood entirely.
        \end{itemize}
    \end{columns}
\end{frame}

\begin{frame}{Energy-Based Models (EBMs)}
    We can explicitly parameterise $p_\theta(x)$ using an energy function $f_\theta \colon \mathbb{R}^d \to \mathbb{R}$ mapped by a neural network yielding
    \begin{align*}
        p_\theta(x) = \frac{\exp(-f_\theta(x))}{Z_\theta} \, , \quad \text{where} \quad Z_\theta = \int \exp(-f_\theta(x)) dx \, .
    \end{align*}
    
    \pause
    Here, $Z_\theta$ is the necessary normalisation constant, widely known as the partition function. 
\end{frame}

\begin{frame}{Training EBMs: The Gradient}
    To minimise the negative log-likelihood $-\frac{1}{N} \sum_{i=1}^N \log p_\theta(x_i)$, we need the gradient $\nabla_\theta \log p_\theta(x)$. Taking the derivative provides
    \begin{align*}
        \nabla_\theta \log p_\theta(x) &= \nabla_\theta \log \left( \frac{\exp(-f_\theta(x))}{Z_\theta} \right) \, , \\
        &= -\nabla_\theta f_\theta(x) - \nabla_\theta \log Z_\theta \, .
    \end{align*}
    
    \onslide<3->{
    The first term is easily computed using standard automatic differentiation on the network output. The second term is vastly more complicated due to the intractable integral in $Z_\theta$.
    }
\end{frame}

\begin{frame}{Training EBMs: The Partition Function}
    We carefully expand the gradient of the log partition function yielding
    \begin{align*}
        \nabla_\theta \log Z_\theta &= \frac{1}{Z_\theta} \nabla_\theta \int \exp(-f_\theta(x)) dx \, , \\
        &= -\int \nabla_\theta f_\theta(x) \frac{\exp(-f_\theta(x))}{Z_\theta} dx \, , \\
        &= -\mathbb{E}_{x \sim p_\theta}[\nabla_\theta f_\theta(x)] \, .
    \end{align*}
    
    \onslide<4->{
    Therefore, the full gradient becomes $-\nabla_\theta f_\theta(x) + \mathbb{E}_{x \sim p_\theta}[\nabla_\theta f_\theta(x)]$.
    
    \begin{alertblock}{Computational Cost}
        Training EBMs is expensive because at every single training iteration, we must sample from the current model distribution $p_\theta$ using MCMC methods to estimate this expectation.
    \end{alertblock}
    }
\end{frame}

\begin{frame}{Implicit Models}
    Implicit models avoid explicit density calculation entirely. They rely on a generator $G_\theta \colon \mathbb{R}^k \to \mathbb{R}^d$ that maps a low-dimensional latent space $z \in \mathbb{R}^k$ to the data space $x \in \mathbb{R}^d$ (where $k \ll d$). 
    
    \pause
    \vspace{0.5em}
    Since the range of the generator $G_\theta$ is confined to a strict $k$-dimensional manifold within the vast $\mathbb{R}^d$ space, the explicit likelihood $p_\theta(x) = 0$ almost everywhere.
\end{frame}

\begin{frame}{Variational Autoencoders (VAEs)}
    To fix the degenerate density problem, VAEs define a noisy observation model $p_\theta(x|z) = \mathcal{N}(x | G_\theta(z), \sigma^2 I)$. The marginal likelihood is now defined everywhere via $p_\theta(x) = \int p_\theta(x|z) p_z(z) dz$, where $p_z(z)$ is a simple prior (e.g., standard normal).
    
    \pause
    \vspace{0.5em}
    Because this integral is intractable, we optimise a lower bound to $\log p_\theta(x)$ and introduce an additional encoder network $e_\phi(z|x) = \mathcal{N}(z | \mu_\phi(x), \Sigma_\phi(x))$.
    
    \vspace{0.5em}
    Here, $\mu_\phi(x)$ and $\Sigma_\phi(x)$ are the mean and covariance of the latent variable, parameterised by the encoder network's weights $\phi$.
\end{frame}

\begin{frame}{The ELBO Objective}
    We train the VAE by maximising the \textbf{Evidence Lower Bound (ELBO)} yielding
    \begin{align*}
        \mathcal{L}_{\text{VAE}} = \mathbb{E}_{z \sim e_\phi(z|x)}[\log p_\theta(x|z)] - D_{KL}(e_\phi(z|x) \,\|\, p_z(z)) \, .
    \end{align*}
    
    \pause
    \begin{itemize}
        \item \textbf{Reconstruction:} The first term trains the decoder $G_\theta$ to accurately reconstruct $x$ from the latent vector $z$.
        \item \textbf{Regularisation:} The KL-divergence forces the encoder's latent space to resemble the standard normal prior $p_z(z)$.
        \item \textbf{Reparameterisation Trick:} Ensures gradients can backpropagate through the stochastic sampling process by expressing the latent variable as $z = \mu_\phi(x) + \Sigma_\phi(x)^{1/2} \epsilon$, where $\epsilon \sim \mathcal{N}(0, I)$.
    \end{itemize}
\end{frame}

\begin{frame}{Generative Adversarial Networks (GANs)}
    Instead of dealing with approximate likelihoods like the ELBO, we can train the generator $G_\theta \colon \mathbb{R}^k \to \mathbb{R}^d$ directly for sampling and bypass density estimation entirely. 
    
    \pause
    \vspace{0.5em}
    We train the generator alongside a discriminator network $D_\phi$, parameterised by weights $\phi$, using an adversarial min-max game yielding
    \begin{align*}
        \mathcal{L}_{\text{GAN}}(\theta, \phi) = -\mathbb{E}_{x \sim p_{\text{data}}}[\log D_\phi(x)] - \mathbb{E}_{z \sim p_z}[\log(1 - D_\phi(G_\theta(z)))] \, .
    \end{align*}
    
    \pause
    \begin{itemize}
        \item The discriminator $D_\phi$ learns to classify true data as $1$ and generated data as $0$.
        \item The generator $G_\theta$ learns to aggressively fool the discriminator into outputting $1$ for generated samples.
    \end{itemize}
\end{frame}

\begin{frame}{Normalising Flows (NFs)}
    In Normalising Flows, we define $G_\theta \colon \mathbb{R}^d \to \mathbb{R}^d$ such that it has the exact same dimensionality in both the latent and image space. 
    
    \pause
    \vspace{0.5em}
    If $G_\theta$ is constructed to be a strict diffeomorphism (invertible and differentiable), we can use the change-of-variables formula to compute the exact analytical likelihood yielding
    \begin{align*}
        p_\theta(x) = p_z(G_\theta^{-1}(x)) \left| \det J_{G_\theta^{-1}}(x) \right| = p_z(z) \left| \det J_{G_\theta}(z) \right|^{-1} \, .
    \end{align*}
    where $p_z(z)$ is a simple base distribution and $J_{G_\theta}(z)$ is the Jacobian matrix of the generator evaluated at $z$.
    
    \pause
    \vspace{0.5em}
    This approach requires highly specialised "Invertible Neural Network" architectures to ensure the Jacobian determinant is computationally tractable.
\end{frame}

\begin{frame}{Summary of Generative Models}
    \begin{table}[htpb]
        \centering
        \begin{tabular}{@{}llll@{}}
            \textbf{Name} & \textbf{Type} & \textbf{Model Structure} & \textbf{Training Method} \\
            \hline
            \textbf{EBM} & Explicit & $p_\theta(x) \propto \exp(-f_\theta(x))$ & Approx. gradient via MCMC \\
            \textbf{NF}  & Explicit & Invertible Network $x = G_\theta(z)$ & Exact Maximum Likelihood \\
            \textbf{VAE} & Implicit & Encoder $e_\phi(z|x)$ / Decoder $G_\theta(z)$ & Lower Bound (ELBO) \\
            \textbf{GAN} & Implicit & Generator $G_\theta(z)$ / Disc. $D_\phi(x)$ & Adversarial min-max game \\
        \end{tabular}
    \end{table}
\end{frame}

\begin{frame}{Generative Models as Priors for Inverse Problems}
    Once a Generative Model (like a VAE or GAN) is trained, the generator $G_\theta \colon \mathbb{R}^k \to \mathbb{R}^n$ acts as an exceptionally powerful data-driven prior.
    
    \pause
    Instead of optimising the image $x$ directly in the high-dimensional space, we \textbf{optimise strictly in the low-dimensional latent space $z$} yielding
    \begin{align*}
        \hat{z} = \argmin_{z \in \mathbb{R}^k} \left\{ \frac{1}{2} \|A G_\theta(z) - y^\delta\|_2^2 + \gamma \|z\|_2^2 \right\} \, .
    \end{align*}
    Here, $\gamma > 0$ is a scalar regularisation parameter. The final reconstruction is simply $\hat{x} = G_\theta(\hat{z})$.
\end{frame}

\begin{frame}{Limitations of Implicit Generators}
    Optimising within the latent space of implicit models carries a significant risk.
    
    \pause
    \begin{alertblock}{The Out-Of-Distribution Hallucination Problem}
        Implicit generators restrict the solution strictly to the pre-trained manifold. If the network is trained on healthy biological tissue, and an unseen pathology (e.g., a rare tumour) appears in the test data $y^\delta$, the generator cannot mathematically represent it. It will aggressively hallucinate a healthy feature instead to minimise the data residual!
    \end{alertblock}
\end{frame}

%================================================================%
\section{Score-Based Models and Diffusion}
%================================================================%

\begin{frame}{Bridging Denoising and Probability}
    We have discussed deterministic denoisers (PnP) and probabilistic distributions. A profound statistical result elegantly connects the two.
    
    When a neural network denoiser $D_\sigma(\cdot)$ is trained using Mean Squared Error (MSE) over a large dataset corrupted by Gau\ss ian noise, it mathematically computes the conditional average (posterior mean).
    
    \pause
    \vspace{0.5em}
    By proving \textbf{Tweedie's Formula}, we can demonstrate that MSE denoisers directly calculate the mathematical score of the data distribution.
\end{frame}

\begin{frame}{Marginal Likelihood as Smoothing}
    Assume additive white Gau\ss ian noise. The noisy-data marginal is
    \begin{align*}
        p_\sigma(y) = \int p(y | x) p(x) \, dx \, .
    \end{align*}
    
    \pause
    Equivalently, $p_\sigma$ is the \textbf{convolution} of the clean prior $p(x)$ with a Gau\ss ian kernel. In short: $p_\sigma$ is a smoothed version of $p(x)$.
    
    \begin{figure}
        \centering
        \includegraphics[width=0.45\textwidth]{pnp/smoothed-marginal-likelihood.png}
    \end{figure}
\end{frame}

\begin{frame}{The Problem with the MMSE Estimator}
    To find the MMSE estimator, we must compute the expected value (mean) of the posterior distribution yielding
    \begin{align*}
        D_\sigma^*(y) = \mathbb{E}[x | y] = \int x \frac{p(y | x) p(x)}{p(y)} \, dx \, .
    \end{align*}
    
    \pause
    \begin{alertblock}{The Fundamental Challenge}
        The true prior distribution $p(x)$ for natural images is an unimaginably complex, high-dimensional manifold. Because we cannot write down a simple mathematical formula for $p(x)$, analytically solving this integral is strictly impossible.
    \end{alertblock}
\end{frame}

\begin{frame}{Tweedie's Formula: The Derivation (1/2)}
    Tweedie's Formula proves that we can bypass the need to know $p(x)$. Taking the gradient of the Gau\ss ian likelihood with respect to the observation $y$ yields
    \begin{align*}
        \nabla_y p(y | x) = -\frac{y - x}{\sigma^2} p(y | x) \, .
    \end{align*}
    
    \pause
    We take the gradient of the marginal distribution $p_\sigma(y)$ with respect to $y$, move the gradient inside the integral, and substitute our likelihood gradient to obtain
    \begin{align*}
        \nabla_y p_\sigma(y) &= \int \nabla_y p(y | x) p(x) \, dx \, , \\
        &= \int \left( -\frac{y - x}{\sigma^2} \right) p(y | x) p(x) \, dx \, .
    \end{align*}
\end{frame}

\begin{frame}{Tweedie's Formula: The Derivation (2/2)}
    Using Bayes' rule, we substitute $p(y | x) p(x) = p(x | y) p_\sigma(y)$ into the integral yielding
    \begin{align*}
        \nabla_y p_\sigma(y) &= -\frac{1}{\sigma^2} \int (y - x) p(x | y) p_\sigma(y) \, dx \, , \\
        &= -\frac{p_\sigma(y)}{\sigma^2} \left( y \int p(x | y) \, dx - \int x \, p(x | y) \, dx \right) \, .
    \end{align*}
    Since $\int p(x | y) dx = 1$ (it is a valid probability distribution) and the second integral exactly forms the posterior mean $\mathbb{E}[x | y]$, we observe
    \begin{align*}
        \nabla_y p_\sigma(y) = -\frac{p_\sigma(y)}{\sigma^2} \Big( y - \mathbb{E}[x | y] \Big) \, .
    \end{align*}
\end{frame}

\begin{frame}{Tweedie's Formula: The Result}
    Dividing both sides by $p_\sigma(y)$ and recognising that $\frac{\nabla p}{p} = \nabla \log p$, we arrive at the profound statistical result known as \textbf{Tweedie's Formula}:
    
    \pause
    \begin{theoremblock}{Tweedie's Formula}
        \begin{align*}
            D_\sigma^*(y) = \mathbb{E}[x | y] = y + \sigma^2 \nabla_y \log p_\sigma(y) \, .
        \end{align*}
    \end{theoremblock}
    
    \pause
    It states that the optimal MMSE denoiser $D_\sigma^*(y)$ depends \emph{only} on the gradient of the log-marginal distribution of the \emph{noisy} data. We never needed to evaluate $p(x)$!
\end{frame}

\begin{frame}{Score-Matching by Denoising (SMD)}
    By rearranging Tweedie's formula, we can express the \textbf{score} (the vector field pointing toward higher probability density) entirely using our denoiser yielding
    \begin{align*}
        \nabla_x \log p_\sigma(x) = \frac{D_\sigma(x) - x}{\sigma^2} \, .
    \end{align*}
    
    \pause
    \textbf{Resolving the Jacobian Critique:} 
    Remember the RED gradient $\nabla \rho(x) = x - D(x)$ that lacked mathematical justification? Tweedie's formula proves that this residual is exactly the negative score!
    
    PnP and RED are \textbf{not} minimising a phantom energy; they are performing explicit gradient ascent on the log-probability of the data manifold \citep{reehorst2018regularization}.
\end{frame}

\begin{frame}{The Need for Sampling}
    Up to this point, our algorithms (like PnP and RED) have been completely deterministic. Given a starting image and a noisy observation, the algorithm climbs the probability landscape until it reaches a single peak or point of consensus.
    
    \pause
    \vspace{0.5em}
    In many inverse problems, the observation is highly ambiguous. Instead of merely finding a single best guess, we want to actively explore all possible solutions by \textbf{sampling} from the posterior distribution $p(x | y^\delta)$.
\end{frame}

\begin{frame}{Langevin Dynamics}
    Once we view our denoiser as a score estimator, we can transition from deterministic optimisation to stochastic sampling.
    
    \pause
    To sample from a target distribution, we simulate physical particle dynamics using \textbf{Langevin Dynamics} \citep{welling2011bayesian} yielding
    \begin{align*}
        x^{k+1} = x^k + \frac{\tau}{2} \nabla_x \log p(x^k) + \sqrt{\tau} z^k \, , \quad z^k \sim \mathcal{N}(0, I) \, .
    \end{align*}
    
    \pause
    \begin{columns}
        \begin{column}{0.6\textwidth}
            \begin{itemize}
                \item \textbf{Gradient Drift:} Climbs towards high-probability regions using the score.
                \item \textbf{Stochastic Noise:} Randomly jostles the image, preventing it from becoming trapped at the MAP peak.
            \end{itemize}
        \end{column}
        \begin{column}{0.4\textwidth}
            \begin{figure}
                \centering
                \includegraphics[width=\textwidth]{pnp/langevin-gradient-drift.png}
            \end{figure}
        \end{column}
    \end{columns}
\end{frame}

\begin{frame}{Decomposing the Posterior Score}
    To solve an inverse problem, we need to sample from the posterior $p(x | y^\delta)$. Applying Bayes' rule and taking the logarithm yields
    \begin{align*}
        \log p(x | y^\delta) = \log p(y^\delta | x) + \log p(x) - \log p(y^\delta) \, .
    \end{align*}
    
    \pause
    Taking the gradient with respect to $x$ drops the evidence term, perfectly decoupling the score into two independent compasses:
    \begin{align*}
        \nabla_x \log p(x | y^\delta) = \underbrace{\nabla_x \log p(y^\delta | x)}_{\text{Likelihood Score}} + \underbrace{\nabla_x \log p(x)}_{\text{Prior Score}} \, .
    \end{align*}
    
    \begin{itemize}
        \item \textbf{Likelihood:} Pushes the image to match physical measurements (e.g., $-\nabla_x \frac{1}{2}\|Ax - y^\delta\|^2$).
        \item \textbf{Prior:} Pushes the image to look natural, given by Tweedie's formula via $\frac{1}{\sigma^2}(D_\sigma(x) - x)$.
    \end{itemize}
\end{frame}

\begin{frame}{The PnP-Langevin Update}
    Substituting these components back into the Langevin dynamics formula provides the final sampling update reading
    \begin{align*}
        x^{k+1} = x^k - \frac{\tau}{2} \nabla_x \mathcal{D}_{y^\delta}(Ax^k) + \frac{\tau}{2\sigma^2} \Big( D_\sigma(x^k) - x^k \Big) + \sqrt{\tau} z^k \, .
    \end{align*}
    
    \pause
    \begin{theoremblock}{A Remarkable Connection}
        If we remove the random noise $\sqrt{\tau} z^k$, this formulation is mathematically identical to the gradient descent update for RED! By injecting controlled noise into PnP/RED, we elegantly transition from classical deterministic solvers to modern stochastic samplers.
    \end{theoremblock}
\end{frame}

\begin{frame}{The Multi-Modal Trap}
    However, running the PnP-Langevin update at a \emph{fixed} noise level $\sigma$ fails practically.
    
    \pause
    \begin{columns}
        \begin{column}{0.5\textwidth}
            \begin{alertblock}{Isolated Modes}
                The true probability landscape of natural images features incredibly sharp, dense peaks separated by vast regions of near-zero probability. A fixed Langevin sampler gets permanently trapped in local modes.
            \end{alertblock}
            
            \vspace{0.5em}
            If we use a large $\sigma$, the image wanders freely but never details. If small, it details but cannot explore.
        \end{column}
        \begin{column}{0.5\textwidth}
            \begin{figure}
                \centering
                \includegraphics[width=0.9\textwidth]{pnp/langevin-annealing.png}
            \end{figure}
        \end{column}
    \end{columns}
\end{frame}

\begin{frame}{Diffusion Models \& Annealed Langevin Dynamics}
    \textbf{Solution:} Introduce a decreasing \textbf{noise schedule} $\sigma_T > \sigma_{T-1} > \dots > \sigma_1 > \sigma_0 \approx 0$.
    
    \vspace{0.5em}
    \pause
    \begin{columns}
        \begin{column}{0.65\textwidth}
            We train a single time-dependent neural network denoiser $D_{\sigma_t}(x)$ across all noise levels. Sampling begins with pure Gau\ss ian noise. At each timestep $t$, we dynamically anneal the noise level yielding
            \begin{align*}
                x^{t-1} &= x^t - \frac{\tau_t}{2} \nabla_x \mathcal{D}_{y^\delta}(A x^t) \\
                &\quad + \frac{\tau_t}{2\sigma_t^2} \left( D_{\sigma_t}(x^t) - x^t \right) + \sqrt{\tau_t} z^t \, .
            \end{align*}
            
            \begin{itemize}
                \item \textbf{High Noise:} Global structural exploration.
                \item \textbf{Low Noise:} Refines fine texture.
            \end{itemize}
        \end{column}
        \begin{column}{0.35\textwidth}
            \begin{figure}
                \centering
                \includegraphics[width=0.9\textwidth]{pnp/noise-schedule.png}
            \end{figure}
        \end{column}
    \end{columns}
\end{frame}

\begin{frame}{Diffusion Posterior Sampling (DPS)}
    When this exact annealed formulation is applied to inverse problems, it is widely known as \textbf{Diffusion Posterior Sampling (DPS)} \citep{chung2023diffusion}.
    
    \vspace{0.5em}
    \pause
    \begin{theoremblock}{The Ultimate Evolution}
        \begin{itemize}
            \item \textbf{PnP / RED:} Deterministic heuristics aiming for a single fixed-point.
            \item \textbf{DPS:} Rigorously samples the full posterior distribution, providing high-quality solutions and quantifying mathematical uncertainty, all while maintaining perfect zero-shot modularity for any linear operator $A$.
        \end{itemize}
    \end{theoremblock}
\end{frame}

%================================================================%
\section{Summary}
%================================================================%

\begin{frame}{Summary}
    Today we covered:
    \begin{enumerate}
        \item \textbf{Modular DL via Plug-and-Play:} Decoupling the forward operator from the network architecture by replacing proximal operators with off-the-shelf CNN denoisers.
        \item \textbf{The Jacobian Critique:} Understanding why explicitly defined functional heuristics like RED fall apart mathematically due to asymmetric Jacobians.
        \item \textbf{Generative Modelling:} Transitioning from deterministic optimisation to explicit and implicit probability density estimation.
        \item \textbf{Score Matching \& Diffusion:} Proving via Tweedie's formula that MSE denoisers calculate the score, enabling posterior sampling through annealed Langevin dynamics (DPS).
    \end{enumerate}
\end{frame}

\begin{frame}{Preview: Thursday}
    We have now explored explicit layer stacking, algorithm unrolling, and modular score matching.
    
    \vspace{0.5em}
    \textbf{Coming up on Thursday:}
    \begin{itemize}
        \item \textbf{Continuous Depth:} Breaking the paradigm of explicit layer-by-layer computation by viewing deep networks as Ordinary Differential Equations (Neural ODEs) \citep{chen2018neural}.
        \item \textbf{Implicit Deep Learning:} How do we build networks with truly infinite depth and $\mathcal{O}(1)$ memory footprints? A deep dive into Deep Equilibrium Models (DEQs) \citep{bai2019deep,gilton2021deep}.
    \end{itemize}
\end{frame}

\begin{frame}[allowframebreaks]{References}
    \bibliographystyle{plainnat}
    \bibliography{references.bib}
\end{frame}

\end{document}